\section{Calibration and library comparisons}\label{app:calibration}

Calibration describes the probability scale on which the evidence estimates depend. The full and direct reduced models are Platt-recalibrated on a separate sample; projection models learn those full-model probabilities without separate outcome recalibration. The production model's test diagnostics are reported in Table~\ref{tab:calibration}.

In the simulation, the unrecalibrated margin has test calibration slope 1.00--1.02 and the fitted Platt slope is 1.00--1.03. After recalibration the test slope is 0.97--1.02, the intercept is $-0.03$ to $0.02$ and decile expected calibration error is at most 0.005. Mean bias against true log odds ranges from $-0.025$ to $0.047$, with regression slope of true on fitted log odds 0.99--1.01. Among declines, predicted default probability averages 0.32--0.33 versus true probability 0.32, and the true-on-fitted log-odds slope is 1.04--1.08. These aggregate agreements coexist with the applicant-level errors in the main text.

\begin{table}[t]\centering\small
\caption{Calibration of the production model on the test sample: ranges over five simulation replicates and ten HELOC splits. The calibration slope and intercept are from a logistic regression of the outcome on the model's log odds (intercept with slope fixed at one). The lower panel compares the simulation model with the true log odds, for all test applicants and for applicants above the cutoff.}
\label{tab:calibration}
\begin{tabular}{@{}lrr@{}}\toprule
& Simulation & HELOC\\\midrule
Platt slope $b$ & 1.00 to 1.03 & 0.81 to 1.18\\
Calibration slope before Platt & 1.00 to 1.02 & 0.91 to 1.08\\
Calibration slope & 0.97 to 1.02 & 0.87 to 1.14\\
Calibration intercept & $-$0.03 to 0.02 & $-$0.19 to 0.08\\
Expected calibration error (deciles) & 0.002 to 0.005 & 0.021 to 0.047\\
Mean predicted probability & 0.087 to 0.092 & 0.484 to 0.528\\
Observed default rate & 0.086 to 0.091 & 0.478 to 0.512\\
\midrule
Bias against true log odds, all & $-$0.025 to 0.047 & --\\
Slope of true on fitted log odds, all & 0.991 to 1.012 & --\\
Bias against true log odds, declined & $-$0.008 to 0.049 & --\\
Slope of true on fitted log odds, declined & 1.039 to 1.076 & --\\
Mean predicted probability, declined & 0.321 to 0.332 & --\\
Mean true probability, declined & 0.317 to 0.323 & --\\
\bottomrule\end{tabular}
\end{table}


\subsection{Gradient-boosting libraries}

LightGBM and CatBoost fitted with the same depth, learning rate and stopping rule reach test AUCs of 0.862 to 0.869 and 0.863 to 0.869. Their direct conditional evidence estimates have RMSE between 0.17 and 0.21 (LightGBM) and between 0.12 and 0.16 (CatBoost) across blocks, comparable to the 0.15 to 0.21 of XGBoost. Against the same applicants and references, the leading candidate reason agrees with the XGBoost model's leading reason for 95 to 97 percent of applicants under LightGBM and 94 to 96 percent under CatBoost, and with the true model's leading reason for 92 to 96 and 93 to 95 percent. Of the XGBoost score-band references, 87 percent pass under the LightGBM model and 78 to 82 percent under the CatBoost model at the same approval rate, so the choice of library changes which records qualify as passing references more than it changes the attribution against a fixed reference set.


\section{Zero-evidence baselines and error decomposition}\label{app:errors}

The zero-evidence baseline has RMSE $\sqrt{\E[W_j^2]}$ for the pointwise target and $\sqrt{\E[(W_j^B)^2]}$ for the class-level target. It measures the error of assigning no conditional evidence to any applicant. Table~\ref{tab:error_decomposition} evaluates Eq.~\eqref{eq:error_difference} on the same combined evaluation records as the main accuracy table. Entries are means of replicate mean squared errors and cross-products, so the identity holds across each row before rounding. Positive cross-products represent cancellation in the evidence difference; correlation alone does not quantify that contribution.

\begin{table}[t]\centering\small
\caption{Error in the evidence difference, averaged over five replicates on the combined evaluation sample. The first two terms are squared full and reduced log-odds errors; subtracting twice their cross-product gives the evidence MSE. Values are in squared log-odds units.}
\label{tab:error_decomposition}
\begin{tabular}{@{}llrrrr@{}}\toprule
Block & Estimator & $\E[e_F^2]$ & $\E[e_{-j}^2]$ & $2\E[e_Fe_{-j}]$ & Evidence MSE\\\midrule
Revolving & Direct & 0.0477 & 0.0405 & 0.0490 & 0.0392\\
 & Projection & 0.0477 & 0.0297 & 0.0490 & 0.0284\\
Delinquency & Direct & 0.0477 & 0.0375 & 0.0412 & 0.0440\\
 & Projection & 0.0477 & 0.0303 & 0.0439 & 0.0341\\
File age & Direct & 0.0477 & 0.0476 & 0.0602 & 0.0351\\
 & Projection & 0.0477 & 0.0361 & 0.0662 & 0.0177\\
Credit seeking & Direct & 0.0477 & 0.0424 & 0.0628 & 0.0273\\
 & Projection & 0.0477 & 0.0360 & 0.0672 & 0.0165\\
Installment & Direct & 0.0477 & 0.0446 & 0.0686 & 0.0236\\
 & Projection & 0.0477 & 0.0411 & 0.0800 & 0.0088\\
Public record & Direct & 0.0477 & 0.0428 & 0.0650 & 0.0255\\
 & Projection & 0.0477 & 0.0401 & 0.0770 & 0.0108\\
\bottomrule\end{tabular}
\end{table}

\begin{table}[t]\centering\small
\caption{Pointwise conditional-evidence RMSE among the 400 sampled declines per replicate, compared with the combined evaluation sample. Entries are means of replicate RMSE in logit units. The decline sample is selected by the original production model and cutoff.}
\label{tab:decline_accuracy}
\resizebox{\linewidth}{!}{\begin{tabular}{@{}lrrrrrr@{}}\toprule
& Revolving & Delinq. & File age & Seeking & Install. & Public rec.\\\midrule
Direct, combined & 0.198 & 0.210 & 0.187 & 0.165 & 0.153 & 0.159\\
Direct, sampled declines & 0.236 & 0.278 & 0.232 & 0.195 & 0.190 & 0.195\\
Projection, combined & 0.168 & 0.185 & 0.133 & 0.128 & 0.093 & 0.104\\
Projection, sampled declines & 0.190 & 0.263 & 0.169 & 0.158 & 0.108 & 0.143\\
\bottomrule\end{tabular}}
\end{table}

